\documentclass[11pt]{letter}
\usepackage{amsmath}
\usepackage[version=3]{mhchem}
\usepackage[margin=1in]{geometry}
\usepackage{hyperref}

\begin{document}

\begin{letter}{Editorial Office \\
Journal of the American Chemical Society / ACS Catalysis \\
American Chemical Society}

\opening{Dear Editor-in-Chief,}

We are pleased to submit our original research manuscript titled \textbf{``Decoupling Electrocatalytic Scaling in Metal-Organic Frameworks via Equivariant Potentials, Micro-Solvation, and Electronic Structure''} for consideration as an Article in your journal.

Electrochemical transformations, including the Hydrogen Evolution Reaction (HER), Oxygen Evolution Reaction (OER), and \ce{CO2} Reduction Reaction (\ce{CO2RR}), are universally constrained by scaling relationships between reactive intermediate binding energies. On conventional solid electrodes, the rigid $\Delta G_{*\text{OOH}} = 0.95\Delta G_{*\text{OH}} + 3.20\,\text{eV}$ relation imposes a thermodynamic overpotential minimum of $\eta^{\text{OER}} \ge 0.37\,\text{V}$, while scaling between $*H$ and $*COOH$ limits selectivity.

In this work, we present an auditable, multi-tier computational protocol coupling screening across 20,372 periodic Metal-Organic Frameworks (MOFs) from the Materials Project / QMOF database with equivariant graph neural network potentials (CHGNet and MACE-MP-0), partial Hessian vibrational analysis (PHVA), explicit pore micro-solvation, and spin-polarized local Density Functional Theory (GPAW, PBE/LCAO). 

Our primary mechanistic finding is that nanoconfined micro-solvation (1--3 explicit \ce{H2O} molecules) within the 1D open channels of Co-MOF-74 selectively stabilizes the hydroperoxo intermediate ($*OOH$, $\Delta E_{\text{solv}} = -2.53\,\text{eV}$) relative to hydroxyl ($*OH$, $-0.65\,\text{eV}$) via a cooperative 6-membered cyclic hydrogen-bonding network. This interaction contracts the scaling gap from $3.20\,\text{eV}$ to $2.74\,\text{eV}$, reducing the OER overpotential from $0.81\,\text{V}$ down to $0.42\,\text{V}$. Furthermore, preferential solvation of the polar carboxyl moiety suppresses parasitic HER and enhances \ce{CO2} reduction selectivity. First-principles GPAW calculations establish that high-lying Co $3d$ bands ($\varepsilon_d - E_F = -1.14\,\text{eV}$) facilitate oxidation, while deeper Cu $3d$ states ($-3.24\,\text{eV}$) maintain near-thermoneutral hydrogen adsorption ($\eta^{\text{HER}} = 0.12\,\text{V}$).

This manuscript has not been published previously and is not under consideration elsewhere. All authors have approved the manuscript and declare no competing financial interests. Complete simulation workflows, relaxed CIF files, and dataset logs are archived in a public repository to ensure transparency and reproducibility.

\closing{Sincerely,}

\noindent
\textbf{Prof. Dr. Luiz Ribeiro Junior} \\
Laboratory for Advanced Computational Materials and Molecular Modeling \\
Email: luiz@materials-sim.org

\end{letter}
\end{document}
